% Options for packages loaded elsewhere
\PassOptionsToPackage{unicode}{hyperref}
\PassOptionsToPackage{hyphens}{url}
\documentclass[
]{article}
\usepackage{xcolor}
\usepackage[margin=1in]{geometry}
\usepackage{amsmath,amssymb}
\setcounter{secnumdepth}{-\maxdimen} % remove section numbering
\usepackage{iftex}
\ifPDFTeX
  \usepackage[T1]{fontenc}
  \usepackage[utf8]{inputenc}
  \usepackage{textcomp} % provide euro and other symbols
\else % if luatex or xetex
  \usepackage{unicode-math} % this also loads fontspec
  \defaultfontfeatures{Scale=MatchLowercase}
  \defaultfontfeatures[\rmfamily]{Ligatures=TeX,Scale=1}
\fi
\usepackage{lmodern}
\ifPDFTeX\else
  % xetex/luatex font selection
\fi
% Use upquote if available, for straight quotes in verbatim environments
\IfFileExists{upquote.sty}{\usepackage{upquote}}{}
\IfFileExists{microtype.sty}{% use microtype if available
  \usepackage[]{microtype}
  \UseMicrotypeSet[protrusion]{basicmath} % disable protrusion for tt fonts
}{}
\makeatletter
\@ifundefined{KOMAClassName}{% if non-KOMA class
  \IfFileExists{parskip.sty}{%
    \usepackage{parskip}
  }{% else
    \setlength{\parindent}{0pt}
    \setlength{\parskip}{6pt plus 2pt minus 1pt}}
}{% if KOMA class
  \KOMAoptions{parskip=half}}
\makeatother
\usepackage{color}
\usepackage{fancyvrb}
\newcommand{\VerbBar}{|}
\newcommand{\VERB}{\Verb[commandchars=\\\{\}]}
\DefineVerbatimEnvironment{Highlighting}{Verbatim}{commandchars=\\\{\}}
% Add ',fontsize=\small' for more characters per line
\usepackage{framed}
\definecolor{shadecolor}{RGB}{248,248,248}
\newenvironment{Shaded}{\begin{snugshade}}{\end{snugshade}}
\newcommand{\AlertTok}[1]{\textcolor[rgb]{0.94,0.16,0.16}{#1}}
\newcommand{\AnnotationTok}[1]{\textcolor[rgb]{0.56,0.35,0.01}{\textbf{\textit{#1}}}}
\newcommand{\AttributeTok}[1]{\textcolor[rgb]{0.13,0.29,0.53}{#1}}
\newcommand{\BaseNTok}[1]{\textcolor[rgb]{0.00,0.00,0.81}{#1}}
\newcommand{\BuiltInTok}[1]{#1}
\newcommand{\CharTok}[1]{\textcolor[rgb]{0.31,0.60,0.02}{#1}}
\newcommand{\CommentTok}[1]{\textcolor[rgb]{0.56,0.35,0.01}{\textit{#1}}}
\newcommand{\CommentVarTok}[1]{\textcolor[rgb]{0.56,0.35,0.01}{\textbf{\textit{#1}}}}
\newcommand{\ConstantTok}[1]{\textcolor[rgb]{0.56,0.35,0.01}{#1}}
\newcommand{\ControlFlowTok}[1]{\textcolor[rgb]{0.13,0.29,0.53}{\textbf{#1}}}
\newcommand{\DataTypeTok}[1]{\textcolor[rgb]{0.13,0.29,0.53}{#1}}
\newcommand{\DecValTok}[1]{\textcolor[rgb]{0.00,0.00,0.81}{#1}}
\newcommand{\DocumentationTok}[1]{\textcolor[rgb]{0.56,0.35,0.01}{\textbf{\textit{#1}}}}
\newcommand{\ErrorTok}[1]{\textcolor[rgb]{0.64,0.00,0.00}{\textbf{#1}}}
\newcommand{\ExtensionTok}[1]{#1}
\newcommand{\FloatTok}[1]{\textcolor[rgb]{0.00,0.00,0.81}{#1}}
\newcommand{\FunctionTok}[1]{\textcolor[rgb]{0.13,0.29,0.53}{\textbf{#1}}}
\newcommand{\ImportTok}[1]{#1}
\newcommand{\InformationTok}[1]{\textcolor[rgb]{0.56,0.35,0.01}{\textbf{\textit{#1}}}}
\newcommand{\KeywordTok}[1]{\textcolor[rgb]{0.13,0.29,0.53}{\textbf{#1}}}
\newcommand{\NormalTok}[1]{#1}
\newcommand{\OperatorTok}[1]{\textcolor[rgb]{0.81,0.36,0.00}{\textbf{#1}}}
\newcommand{\OtherTok}[1]{\textcolor[rgb]{0.56,0.35,0.01}{#1}}
\newcommand{\PreprocessorTok}[1]{\textcolor[rgb]{0.56,0.35,0.01}{\textit{#1}}}
\newcommand{\RegionMarkerTok}[1]{#1}
\newcommand{\SpecialCharTok}[1]{\textcolor[rgb]{0.81,0.36,0.00}{\textbf{#1}}}
\newcommand{\SpecialStringTok}[1]{\textcolor[rgb]{0.31,0.60,0.02}{#1}}
\newcommand{\StringTok}[1]{\textcolor[rgb]{0.31,0.60,0.02}{#1}}
\newcommand{\VariableTok}[1]{\textcolor[rgb]{0.00,0.00,0.00}{#1}}
\newcommand{\VerbatimStringTok}[1]{\textcolor[rgb]{0.31,0.60,0.02}{#1}}
\newcommand{\WarningTok}[1]{\textcolor[rgb]{0.56,0.35,0.01}{\textbf{\textit{#1}}}}
\usepackage{graphicx}
\makeatletter
\newsavebox\pandoc@box
\newcommand*\pandocbounded[1]{% scales image to fit in text height/width
  \sbox\pandoc@box{#1}%
  \Gscale@div\@tempa{\textheight}{\dimexpr\ht\pandoc@box+\dp\pandoc@box\relax}%
  \Gscale@div\@tempb{\linewidth}{\wd\pandoc@box}%
  \ifdim\@tempb\p@<\@tempa\p@\let\@tempa\@tempb\fi% select the smaller of both
  \ifdim\@tempa\p@<\p@\scalebox{\@tempa}{\usebox\pandoc@box}%
  \else\usebox{\pandoc@box}%
  \fi%
}
% Set default figure placement to htbp
\def\fps@figure{htbp}
\makeatother
\setlength{\emergencystretch}{3em} % prevent overfull lines
\providecommand{\tightlist}{%
  \setlength{\itemsep}{0pt}\setlength{\parskip}{0pt}}
\usepackage{bookmark}
\IfFileExists{xurl.sty}{\usepackage{xurl}}{} % add URL line breaks if available
\urlstyle{same}
\hypersetup{
  pdftitle={assignment\_3},
  pdfauthor={Emi},
  hidelinks,
  pdfcreator={LaTeX via pandoc}}

\title{assignment\_3}
\author{Emi}
\date{2026-06-23}

\begin{document}
\maketitle

\section{\texorpdfstring{\textbf{Exercise 1. Corruption and human
development}}{Exercise 1. Corruption and human development}}\label{exercise-1.-corruption-and-human-development}

\begin{center}\rule{0.5\linewidth}{0.5pt}\end{center}

This exercise explores a dataset containing the human development index
(\texttt{HDI}) and corruption perception index (\texttt{CPI}) of 173
countries across 6 different regions around the world: Americas, Asia
Pacific, Eastern Europe and Central Asia (\texttt{East\ EU\ Cemt}),
Western Europe (\texttt{EU\ W.\ Europe}), Middle East and North Africa
and Noth Africa (\texttt{MENA}), and Sub-Saharan Africa (\texttt{SSA}).
(Note: the larger \texttt{CPI} is, the less corrupted the country is
perceived to be.)

First, we load the data using the following code.

\begin{Shaded}
\begin{Highlighting}[]
\NormalTok{economist\_data }\OtherTok{\textless{}{-}} \FunctionTok{read.csv}\NormalTok{(}\StringTok{"https://raw.githubusercontent.com/nt246/NTRES6940{-}data{-}science/master/datasets/EconomistData.csv"}\NormalTok{)}
\end{Highlighting}
\end{Shaded}

\textbf{1.1 Show the first few rows of \texttt{economist\_data}.}

\begin{Shaded}
\begin{Highlighting}[]
\FunctionTok{head}\NormalTok{(economist\_data)}
\end{Highlighting}
\end{Shaded}

\begin{verbatim}
##   X     Country HDI.Rank   HDI CPI            Region
## 1 1 Afghanistan      172 0.398 1.5      Asia Pacific
## 2 2     Albania       70 0.739 3.1 East EU Cemt Asia
## 3 3     Algeria       96 0.698 2.9              MENA
## 4 4      Angola      148 0.486 2.0               SSA
## 5 5   Argentina       45 0.797 3.0          Americas
## 6 6     Armenia       86 0.716 2.6 East EU Cemt Asia
\end{verbatim}

\textbf{1.2 Expore the relationship between human development index
(\texttt{HDI}) and corruption perception index (\texttt{CPI}) with a
scatter plot as the following.}

\begin{Shaded}
\begin{Highlighting}[]
\DocumentationTok{\#\# Write your code here}
\end{Highlighting}
\end{Shaded}

\textbf{1.3 Make the color of all points in the previous plot red}

\begin{Shaded}
\begin{Highlighting}[]
\DocumentationTok{\#\# Write your code here}
\end{Highlighting}
\end{Shaded}

\textbf{1.4 Color the points in the previous plot according to the
\texttt{Region} variable, and set the size of points to 2.}

\begin{Shaded}
\begin{Highlighting}[]
\DocumentationTok{\#\# Write your code here}
\end{Highlighting}
\end{Shaded}

\textbf{1.5 Set the size of the points proportional to
\texttt{HDI.Rank}}

\begin{Shaded}
\begin{Highlighting}[]
\DocumentationTok{\#\# Write your code here}
\end{Highlighting}
\end{Shaded}

\textbf{1.6 Fit a smoothing line to all the data points in the scatter
plot from Excercise 1.4}

\begin{Shaded}
\begin{Highlighting}[]
\DocumentationTok{\#\# Write your code here}
\end{Highlighting}
\end{Shaded}

\textbf{1.7 Fit a separate straight line for each region instead, and
turn off the confidence interval.}

\begin{Shaded}
\begin{Highlighting}[]
\DocumentationTok{\#\# Write your code here}
\end{Highlighting}
\end{Shaded}

\textbf{1.8 Building on top of the previous plot, show each
\texttt{Region} in a different facet.}

\begin{Shaded}
\begin{Highlighting}[]
\DocumentationTok{\#\# Write your code here}
\end{Highlighting}
\end{Shaded}

\textbf{1.9 Show the distribution of \texttt{HDI} in each region using
density plot. Set the transparency to 0.5}

\begin{Shaded}
\begin{Highlighting}[]
\DocumentationTok{\#\# Write your code here}
\end{Highlighting}
\end{Shaded}

\textbf{1.10 Show the distribution of \texttt{HDI} in each region using
histogram and facetting.}

\begin{Shaded}
\begin{Highlighting}[]
\DocumentationTok{\#\# Write your code here}
\end{Highlighting}
\end{Shaded}

\textbf{1.11 Show the distribution of \texttt{HDI} in each region using
a box plot. Set the transparency of these boxes to 0.5. Also show data
points for each country in the same plot. (Hint: \texttt{geom\_jitter()}
or \texttt{position\_jitter()} might be useful.)}

\begin{Shaded}
\begin{Highlighting}[]
\DocumentationTok{\#\# Write your code here}
\end{Highlighting}
\end{Shaded}

\textbf{1.12 Show the count of countries in each region using a bar
plot.}

\begin{Shaded}
\begin{Highlighting}[]
\DocumentationTok{\#\# Write your code here}
\end{Highlighting}
\end{Shaded}

\textbf{1.13 You have now created a variety of different plots of the
same dataset. Which of your plots do you think are the most informative?
Describe briefly the major trends that you see in the data.}

\begin{Shaded}
\begin{Highlighting}[]
\DocumentationTok{\#\# Write your code here}
\end{Highlighting}
\end{Shaded}

\section{\texorpdfstring{\textbf{Exercise 2. Unemployment in the US
1967-2015}}{Exercise 2. Unemployment in the US 1967-2015}}\label{exercise-2.-unemployment-in-the-us-1967-2015}

\begin{center}\rule{0.5\linewidth}{0.5pt}\end{center}

This excercise uses the dataset \texttt{economics} from the ggplot2
package. It was produced from US economic time series data available
from \url{http://research.stlouisfed.org/fred2}. It descibes the number
of unemployed persons (\texttt{unemploy}), among other variables, in the
US from 1967 to 2015.

\begin{Shaded}
\begin{Highlighting}[]
\DocumentationTok{\#\# Write your code here}
\end{Highlighting}
\end{Shaded}

\textbf{2.1 Plot the trend in number of unemployed persons
(\texttt{unemploy}) though time using the economics dataset shown above.
And for this question only, hide your code and only show the plot.}

\begin{Shaded}
\begin{Highlighting}[]
\DocumentationTok{\#\# Write your code here}
\end{Highlighting}
\end{Shaded}

\textbf{2.2 Edit the plot title and axis labels of the previous plot
appropriately. Make y axis start from 0. Change the background theme to
what is shown below. (Hint: search for help online if needed)}

\begin{Shaded}
\begin{Highlighting}[]
\DocumentationTok{\#\# Write your code here}
\end{Highlighting}
\end{Shaded}


\end{document}
